\section{Why the logarithm is a Lie series}\label{sec:hopf}
We prove the primitive-element statement used in the source article rather than treating it as a black box. The proof simultaneously constructs the projection that produces Dynkin's formula.

\subsection{The coproduct and group-like elements}
Define the continuous algebra homomorphism
\begin{equation}\label{eq:coproduct}
 \Delta:\widehat\TT\longrightarrow\widehat\TT\ot\widehat\TT,
 \qquad
 \Delta X=X\otimes1+1\otimes X,\quad
 \Delta Y=Y\otimes1+1\otimes Y.
\end{equation}
Here the target is the \emph{completed} tensor product
$\prod_{p,q\geq0}\TT_p\otimes\TT_q$, with the total-degree interpretation of multiplication and infinite sums. An ordinary algebraic tensor product is not sufficient for arbitrary formal series. Explicitly, if $w=a_1\cdots a_n$, then
\begin{equation}\label{eq:unshuffle}
 \Delta w=\sum_{I\subseteq\{1,\ldots,n\}}a_I\otimes a_{I^c},
\end{equation}
where $a_I$ lists the selected letters in their original order. The counit $\eps$ selects the constant coefficient. Coassociativity and the counit equations can be checked on $X,Y$ and then extended by multiplicativity and continuity.

An element $P$ is \emph{primitive} when
$\Delta P=P\otimes1+1\otimes P$.
An element $g\in1+F^1\widehat\TT$ is \emph{group-like} when
$\Delta g=g\otimes g$.
The normalization of the constant term is part of this definition.

\begin{proposition}\label{prop:group-like}
For $P\in F^1\widehat\TT$, $P$ is primitive if and only if $e^P$ is group-like. Group-like elements form a group under multiplication.
\end{proposition}
\begin{proof}
Continuity and multiplicativity give $\Delta(e^P)=e^{\Delta P}$. If $P$ is primitive, the two tensor factors commute, so
\[
 e^{\Delta P}=e^{P\otimes1}e^{1\otimes P}=e^P\otimes e^P.
\]
Conversely, if $g=e^P$ is group-like, then
\begin{align*}
 \Delta P&=\log(\Delta g)
 =\log\bigl((g\otimes1)(1\otimes g)\bigr)\\
 &=\log(g\otimes1)+\log(1\otimes g)
 =P\otimes1+1\otimes P.
\end{align*}
The commuting-logarithm rule applies because both factors have constant term one. Finally,
$\Delta(gh)=(g\otimes g)(h\otimes h)=gh\otimes gh$,
and $\Delta(g^{-1})=(\Delta g)^{-1}=g^{-1}\otimes g^{-1}$.
\end{proof}

\subsection{An explicit primitive-element projection}
For a word $w$ of length $n$, define
\[
 S(w)=(-1)^n\rev(w),\qquad N(w)=n w,
 \qquad S(1)=1,\quad N(1)=0.
\]
The map $S$ is an antihomomorphism. For endomorphisms $f,g$ of $\TT$, their convolution is
$ f*g=m(f\otimes g)\Delta$, where $m$ is multiplication.

\begin{lemma}[Antipode and right Dynkin operator]\label{lem:dynkin-operator}
The map $S$ satisfies $\Id*S=S*\Id=\eta\eps$, where $\eta$ inserts scalars. The convolution
\begin{equation}\label{eq:dynkin-operator}
 D=N*S
\end{equation}
satisfies $D(w)=R(w)$ for every nonempty word.
\end{lemma}
\begin{proof}
Use the finite notation $\Delta w=\sum u\otimes v$ and put
$K(w)=\sum uS(v)$. If $a$ is a letter, then
$\Delta(aw)=\sum au\otimes v+u\otimes av$ and $S(av)=-S(v)a$.
Consequently
\[
 K(aw)=aK(w)-K(w)a=0.
\]
Since $K(1)=1$, this proves $\Id*S=\eta\eps$. The corresponding calculation using a final letter proves the other antipode identity (or one can use uniqueness of the convolution inverse in the connected graded algebra).

Next, $N(au)=aN(u)+au$. Therefore
\begin{align*}
 D(aw)
 &=\sum N(au)S(v)+N(u)S(av)\\
 &=aD(w)-D(w)a+aK(w).
\end{align*}
For $w\ne1$, the last term vanishes. For $w=1$, it equals $a$, so $D(a)=a$. Induction now gives $D(aw)=[a,D(w)]$, which is exactly right nesting.
\end{proof}

\begin{theorem}[Primitive elements and Dynkin--Specht--Wever]\label{thm:friedrichs}
A homogeneous $P\in\TT_n$, $n>0$, is primitive if and only if it belongs to $\LL(X,Y)_n$. For such a $P$,
\begin{equation}\label{eq:dsw}
 D(P)=nP,\qquad
 P=\frac1n\sum_{|w|=n}\coeff{w}P\,R(w).
\end{equation}
Thus $D/n$ is an idempotent projection of $\TT_n$ onto $\LL(X,Y)_n$, and the primitive elements of $\widehat\TT$ are exactly $\widehat\LL(X,Y)$.
\end{theorem}
\begin{proof}
If $P$ is homogeneous and primitive, then
\[
 D(P)=m(N\otimes S)(P\otimes1+1\otimes P)=nP.
\]
By \cref{lem:dynkin-operator}, $D(P)$ is a linear combination of nested commutators. Division by $n$ is permitted in characteristic zero, proving $P\in\LL$ and the displayed formula.

Conversely, $X,Y$ are primitive, and the bracket of primitive elements is primitive: the mixed terms in $[\Delta U,\Delta V]$ cancel because the tensor factors commute. Hence every Lie polynomial is primitive. Since $D$ always takes values in Lie polynomials, applying $D$ once more gives $D^2=nD$ on degree $n$. Finally, $\Delta$ preserves total degree, so a completed series is primitive exactly when each homogeneous component is primitive. Its degree-zero component must be zero, since a primitive scalar $c$ would satisfy $c=2c$.
\end{proof}

\begin{theorem}[Formal BCH theorem]\label{thm:formal-bch}
There is a unique $Z\in F^1\widehat\TT$ satisfying $e^Z=e^Xe^Y$. It belongs to $\widehat\LL(X,Y)$, has rational coefficients, and its degree-one term is $X+Y$.
\end{theorem}
\begin{proof}
Uniqueness and existence in the completed associative algebra follow from \cref{lem:exp-log}. Since $X,Y$ are primitive, their exponentials and the product $e^Xe^Y$ are group-like. Its logarithm is primitive by \cref{prop:group-like}, and hence a Lie series by \cref{thm:friedrichs}. Rationality follows from the factorial and logarithmic coefficients. Taking degree one gives $Z_1=X+Y$.
\end{proof}
This proves all the structural assertions about primitive and group-like elements needed for BCH. The Hopf-algebra interpretation is compatible with universal enveloping algebras, as explained and proved in \cref{sec:pbw}.
